В этой работе реализован вариант FMPE с последовательным уточнением
априорного распределения: на каждом раунде оно усекалось по
posterior-выборкам, полученным на предыдущем.

Эксперименты не показали, что дополнительные раунды систематически
повышают качество; в некоторых конфигурациях работа алгоритма становилась
нестабильной. На гауссовской задаче при одном раунде средний C2ST SNPE-C
был ближе к идеальному значению, чем у выбранной конфигурации SFMPE.

На SIR и StochVol соотношение методов по ATC менялось вместе с бюджетом
симуляций. Эти результаты следует трактовать с учётом малого числа
целевых наблюдений и того, что часть posterior-сэмплов оказывалась вне
допустимой области параметров. Поэтому имеющиеся эксперименты не дают
оснований говорить о безусловном превосходстве предложенного метода.
